\section{Introduction}
\label{sec-intro}

Ontologies such as the Gene Ontology~\cite{ashburner2000} and the schemas of
enterprise knowledge graphs are stored as sets of triples and queried in
three ways.
A \emph{pattern} query asks for the triples that match a template and is the
base case of every RDF store~\cite{weiss2008,neumann2010}.
A \emph{structural} query asks for the instances at or below a class, and
needs the class hierarchy.
A \emph{semantic} query asks for the terms closest in meaning to a given
term, which needs a similarity measure over embeddings.
The three meet in ``the terms most similar to \emph{apoptotic process} that
lie under \emph{cell death}''.

Existing systems leave four problems open.
\textbf{C1. Hierarchy costs space or time.}
RDFS entailment is either \emph{materialized}, which stores every inferred
triple~\cite{rdfox}, or \emph{rewritten} into the query, which scans once per
subclass.
On the Gene Ontology with 200,000 typed entities, materializing multiplies
the triples by 8.1, and rewriting reads 48 times more pages than a store that
knows the hierarchy.
\textbf{C2. A sorted index does not hold meaning.}
It is tempting to believe that neighbors in meaning are neighbors in the
order; we prove they cannot all be (Section~\ref{sec-sem}), so similarity is
left to a vector database beside the triple store, with its own storage and
no view of the hierarchy.
\textbf{C3. A class filter is correlated with the vectors.}
Filtered vector search is costed by assuming the filter is independent of the
similarity; in an ontology the neighbors of a term lie under its own
ancestors, and a post-filtered index can find 7\% of the right answers.
\textbf{C4. A paged store must scale across threads.}
A classic buffer pool serializes every page access on one latch, and a classic
B+ tree write latches the root.
In our first version, lookups per second fell from 330 thousand with one
thread to 180 thousand with eight.

\begin{example}
\label{exa-motivation}
Let $q$ be \emph{apoptotic process} in the Gene Ontology and $c$ the class
\emph{cell death}.
The query asks for the five terms nearest to $q$ among the 97 classes at or
below $c$ in an 834-class slice.
A triple store cannot rank by meaning; a vector index returns neighbors from
anywhere, and a post-filter keeps those under $c$ only if enough are.
\sys finds the 97 classes with one range scan of the B+ tree that holds the
triples, scores them exactly, and returns \emph{apoptotic process involved in
development} (cosine 0.865) first, in under a millisecond in a web page.
\end{example}

\stitle{Contributions}.
This paper introduces \sys, an open-source ontology database in C++20
(Figure~\ref{fig-arch}), with five technical contributions.
\begin{enumerate}
\item \textbf{A concurrent B+ tree triple store} (Sections~\ref{sec-store}
  and~\ref{sec-conc}): the algorithms for search, insert, and delete, a bound
  on the pages a range scan reads, and a latch-free hit path and an optimistic
  writer descent that scale to eight threads.
\item \textbf{Hierarchies as intervals} (Section~\ref{sec-labels}): a postorder
  labelling of the class DAG kept in the same B+ tree, so that ``everything
  under $c$'' is a few range scans and no entailed triple is stored.
\item \textbf{An entailment-aware plan rewrite} (Section~\ref{sec-rewrite}):
  RDFS subclass entailment of \texttt{rdf:type} patterns as interval scans the
  optimizer costs like any other, with the answers of materialization, prepared
  16 times faster and in 8 times less space.
\item \textbf{Vector indexes and filtered search} (Section~\ref{sec-sem}): a
  proof that no sorted order keeps neighbors close, IVF posting lists stored as
  key ranges of the triple tree, and the finding that a class filter is
  correlated with similarity, with the rule that chooses a strategy.
\item \textbf{A build for the browser} (Section~\ref{sec-system}): the whole
  engine as a 0.85\,MB WebAssembly module.
\end{enumerate}
We evaluate \sys against Oxigraph~\cite{oxigraph}, Apache Jena
TDB2~\cite{jenatdb}, and SQLite on LUBM~\cite{lubm} and the Gene Ontology
(Section~\ref{sec-exp}, Table~\ref{tab-glance}).
The code, data generators, and benchmarks are at the artifact URL above.
